%======================================================================
% Table T4 — Wave types and governing characteristics
% Section: §3, §4
%======================================================================
\begin{table}[htbp]
\centering
\small
\caption{Wave types and representative characteristics across continuum models.}
\label{tab:T4-waves}
\begin{tabularx}{\textwidth}{@{}p{2.4cm}XXX@{}}
\toprule
\textbf{Wave type} & \textbf{Polarization or field pattern} & \textbf{Classical reference feature} & \textbf{Model-specific interpretation} \\
\midrule
Longitudinal (P) & Parallel to propagation & $c_p=\sqrt{(\lambda+2\mu)/\rho}$ for isotropic elasticity & Generalized mode labels and speeds follow the selected constitutive equations \\
Transverse (S) & Perpendicular to propagation & $c_s=\sqrt{\mu/\rho}$ for isotropic elasticity & Coupling can alter the modal structure \\
Thermal & Temperature and heat-flux fields & Diffusion in the Fourier model & Relaxation, phase-lag, and memory formulations have model-dependent propagation behavior \\
Coupled modes & Mixed mechanical, thermal, or electrical fields & Eigenfrequencies of the coupled system & Branches follow the coupled field equations and boundary conditions \\
Surface waves & Amplitude decays with depth from the surface & Rayleigh-type branch in an elastic half-space & Constitutive laws and boundary conditions determine the surface-wave spectrum \\
Evanescent modes & Spatially decaying amplitude & Imaginary part of the wave number & Near-field contributions can enter interface and scattering solutions \\
\bottomrule
\end{tabularx}
\end{table}
